\section{Introduction}
\label{sec:intro}

Ligand-based activity prediction rests on an assumption that is rarely stated
explicitly: that the chemical features which determine whether a compound is
active are carried by the organic scaffold, and therefore that a model trained
on one set of compounds will rank a new set of compounds sensibly even if the
scaffolds differ. This assumption is the reason scaffold splits
\citep{bemis1996murcko} are treated as the honest evaluation for molecular
property models, and the reason transfer across organic chemical space is
considered a solved-enough problem to build benchmarks on
\citep{moleculenet2018}.

For organometallic and coordination compounds the assumption is weaker in a
specific, testable way. Two complexes can share an identical ligand set and
differ only in the coordinated metal, and that single change can move the
measured anticancer activity by orders of magnitude. The metal is not a
substituent; it sets the geometry, the accessible oxidation states, the
hydrolysis rate, and the DNA-binding mode. If a ligand-based model has learned
``which scaffolds are toxic'', it should transfer across metals. If it has
learned ``which metal--ligand combinations are toxic'', it should not.

We test this directly. Using the MetalCytoToxDB cohort of precious- and
platinum-group anticancer complexes \citep{krasnov2026metalcytotoxdb}, we
evaluate the \emph{same} model and the \emph{same} scorer under three regimes:
within-metal 5-fold cross-validation, a same-metal Bemis--Murcko scaffold-split
control, and cross-metal leave-one-metal-out (LOMO) transfer. The within- and
cross-metal regimes differ in more than the held-out metal --- their training
sets, sample sizes and label priors all change --- so we use the scaffold
control to measure how much of the within-to-cross drop is ordinary
novel-scaffold generalisation and how much is specific to the metal axis.
Neither the ligand-based nor the transition-metal literature reports this
cross-metal measurement (\S\ref{sec:related}), so we treat it as the primary
object of study rather than an incidental split.

Our contributions are:

\begin{enumerate}
\item \textbf{A measured collapse, its decomposition, and its robustness.}
Within-metal ROC-AUC is $0.70$--$0.95$ for Ru, Ir, Rh and Re (Os's within-metal
signal does not resolve), but cross-metal LOMO AUC falls to $0.44$--$0.62$ for
the same model. A same-metal Bemis--Murcko scaffold-split control costs
$0.10$--$0.12$ on the two metals where it resolves, so roughly half of the
collapse is ordinary novel-scaffold generalisation and the metal axis is the
largest single remaining shift, not the only one. The collapse is not an
artefact of the scorer: it reproduces under an independent message-passing
backbone on the same cohort and split, which additionally recovers Re
($0.737$, Holm $p=0.010$) where the fingerprint probe scores $0.443$ --- so we
report Re as representation-dependent. Under a Holm-corrected label-permutation
null the fingerprint scorer leaves only Ru significantly above chance, and the
backbone leaves only Re and Ru significant: Ir's below-chance reading is
significant only uncorrected ($p=0.038$) and does not survive the correction. We argue metal
identity is a distribution shift with no routine organic analogue, and that it
should be reported as a benchmark axis in its own right.
\MEASURED. \S\ref{sec:results}.

\item \textbf{No significant pretraining transfer, measured on two metals.}
A graph encoder pretrained on the tmQM transition-metal corpus
\citep{tmqm2020} does not raise the activity head on a genuine scaffold split.
On Ru the paired delta is $-0.0094$ (t-based 95\% CI $[-0.0406,+0.0218]$,
$t(4)=-0.84$, 2/5 seeds positive); on Ir it is a suggestive $+0.0486$ that does
not clear the paired $t$-test ($3/5$ seeds positive, one-sided $p=0.068$, CI
includes zero). An earlier $+0.0181$ was measured on a
degenerate, molecule-level splitter and does not survive the fix. Both arms
print the tensor-load count, so a silent no-op is excluded; the null is a null
of the effect, not of the measurement.
\MEASURED. \S\ref{sec:results}.

\item \textbf{A within-metal win and a cross-cohort honest negative.} On the
source paper's own within-metal axis our fingerprint-plus-XGBoost scorer exceeds
the published $0.73$ under both tie rules and every seed
($0.8237$--$0.8597$); Ru is at parity,
decided by the tie rule; and we report the first within-metal values for Rh, Os
and Re, which the source paper does not cover. Under their DOI-split 10-fold
cross-validation our model scores $0.7717 \pm 0.0380$ against the published
LightGBM $0.810 \pm 0.024$: $0.0383$ below, with 3 of 10 folds at or above
$0.81$, on a cohort $25\times$ wider. We report both the win and the gap rather
than burying either. \MEASURED. \S\ref{sec:results}.

\item \textbf{A method, and a clean negative for the obvious alternative.}
Static metal conditioning --- a metal token, a metal-descriptor embedding, or a
FiLM scale-and-shift --- does not repair the collapse: all three leave the
cross-metal mean at or below the unconditioned baseline (FiLM $-0.0418$, Holm
$p=0.984$). What does work is few-shot adaptation. Fine-tuning on $K$ labelled
ligands from the held-out metal recovers AUC monotonically, by $+0.058$ at
$K{=}20$ and $+0.091$ at $K{=}40$ (Holm $p<0.004$, ten seeds). The collapse is
therefore not a missing-input problem but a target-distribution problem, and a
handful of target-metal labels is the cheap fix. \MEASURED.
\S\ref{sec:results:method}.
\end{enumerate}

The paper is organised as follows. \S\ref{sec:related} positions the work
against ligand-based QSAR and transition-metal machine learning.
\S\ref{sec:method} defines the cohort, the two evaluation regimes, and the
model. \S\ref{sec:results} reports the four claims.
\S\ref{sec:ablation} gives the sensitivity analyses that bound the conclusions
(tie rule, deduplication, pretraining A/B). \S\ref{sec:limitations} states what
the study does \emph{not} establish, including a cohort-truncation effect we
found in a related protocol and the small per-metal sample sizes that make part
of the result unresolvable.
